\chapter{LU Factorization}

In the last chapter you learned to solve $\mat{A}\vec{x} = \vec{b}$ by elimination:
sweep down the matrix, knocking out the entries below each pivot, until what remains is
upper triangular and easy to back-solve. It works. But there is a quiet tragedy hidden
inside that procedure. Every time someone hands you a new right-hand side $\vec{b}$, you
seem to have to start the whole sweep again from the top.

This chapter fixes that. The key idea is simple once you see it: \emph{the work of
elimination can be bottled up and reused}. All the row
operations you performed -- every ``subtract $3$ times row $1$ from row $2$'' -- can be
recorded once, packaged into a single matrix, and replayed instantly on any new $\vec{b}$.
That packaging is the LU factorization. We will discover that elimination is secretly
factoring your matrix into two triangular pieces,
\[
  \mat{A} = \mat{L}\,\mat{U},
\]
a \textbf{l}ower triangular $\mat{L}$ times an \textbf{u}pper triangular $\mat{U}$. Once
you have $\mat{L}$ and $\mat{U}$, solving a system is two fast triangular slides instead of
a full elimination. This is not a minor convenience. It is the workhorse that solves linear
systems inside almost every piece of scientific and engineering software on Earth.

We will build it geometrically (triangular matrices are the simplest distortions of space),
then algebraically (the multipliers \emph{are} the entries of $\mat{L}$), then handle the
real-world wrinkle (row exchanges, giving $\mat{P}\mat{A} = \mat{L}\mat{U}$), and finally
count the cost so you understand \emph{why} factoring once and reusing beats recomputing.

\section{The picture: triangular machines are easy machines}
\label{ch09:triangular-picture}

Before any formulas, hold onto a feeling for what triangular matrices \emph{do} to space.

A matrix is a machine that takes a vector in and pushes a vector out. An \emph{upper
triangular} matrix
\[
  \mat{U} = \begin{bmatrix} 2 & 1 \\ 0 & 3 \end{bmatrix}
\]
has all its action ``above the diagonal.'' Look at what it does to the standard basis. The
first basis vector $\colvec{1\\0}$ goes to $\colvec{2\\0}$, staying flat on the
$x$-axis -- the $x$-axis is left as a line through the origin, just stretched. The second
basis vector $\colvec{0\\1}$ goes to $\colvec{1\\3}$, which leans to the right. Upper
triangular means: \emph{the first axis stays put as a direction; later axes are allowed to
spill into the earlier ones.}

A \emph{lower triangular} matrix
\[
  \mat{L} = \begin{bmatrix} 1 & 0 \\ 3 & 1 \end{bmatrix}
\]
is the mirror image. Here $\colvec{1\\0}$ goes to $\colvec{1\\3}$ (the first axis spills
\emph{upward} into the second) while $\colvec{0\\1}$ stays put as $\colvec{0\\1}$. This particular $\mat{L}$, with $1$'s on the diagonal, is a \textbf{shear}: it slides the plane
vertically, keeping the vertical direction fixed. We met shears as one of the basic
deformations a matrix can apply; now they take on a starring role.

\begin{figure}[H]
\centering
\begin{tikzpicture}[scale=1.05,>=Stealth]
  % left: lower-triangular shear acting on unit square
  \begin{scope}[xshift=0cm]
    \draw[->,gray!60] (-0.3,0)--(2.6,0);
    \draw[->,gray!60] (0,-0.3)--(0,3.4);
    % original unit square
    \draw[vecblue,thick] (0,0)--(1,0)--(1,1)--(0,1)--cycle;
    \node[vecblue] at (0.5,-0.35) {\small unit square};
    % sheared image (L = [[1,0],[3,... ]] too steep; use modest shear [[1,0],[1,1]])
    \draw[vecred,thick] (0,0)--(1,1)--(1,2)--(0,1)--cycle;
    \draw[->,vecred,thick] (0.5,0.5)--(0.95,1.45);
    \node[vecred] at (1.7,1.9) {\small image};
    \node at (1.0,-1.0) {\small lower triangular: shear};
  \end{scope}
  % right: upper-triangular acting on unit square
  \begin{scope}[xshift=6.5cm]
    \draw[->,gray!60] (-0.3,0)--(3.0,0);
    \draw[->,gray!60] (0,-0.3)--(0,3.4);
    \draw[vecblue,thick] (0,0)--(1,0)--(1,1)--(0,1)--cycle;
    % U = [[2,1],[0,1.5]] image of square: corners (0,0),(2,0),(3,1.5),(1,1.5)
    \draw[vecred,thick] (0,0)--(2,0)--(3,1.5)--(1,1.5)--cycle;
    \node at (1.5,-1.0) {\small upper triangular: stretch + lean};
  \end{scope}
\end{tikzpicture}
\caption{Triangular matrices are the simplest non-trivial deformations. A lower triangular
matrix with unit diagonal is a pure shear (left): one axis stays fixed, the other slides
along it. An upper triangular matrix stretches the first axis and lets later axes lean into
it (right). Factoring $\mat{A}=\mat{L}\mat{U}$ writes a general distortion as a stretch-and-lean
followed by a shear.}
\label{ch09:fig-triangular}
\end{figure}

\begin{intuitionbox}
LU factorization says: \emph{any} (nice enough) square matrix's action can be done in two
simple stages. First apply $\mat{U}$ -- stretch the space and let later coordinates lean
into earlier ones. Then apply $\mat{L}$ -- a shear that slides earlier coordinates into
later ones. Two triangular slides, composed, can reproduce any general matrix. Elimination
is the act of \emph{discovering} those two slides.
\end{intuitionbox}

Why do we care that the pieces are triangular? Because triangular systems are the easiest
systems in the world to solve. If $\mat{U}$ is upper triangular, the bottom equation
involves only the last unknown, so you read it off directly; the second-to-last equation
then involves only the last two, and you already know one of them; and so on \emph{up} the
list. That is \textbf{back substitution}. If $\mat{L}$ is lower triangular, the same trick
runs top-to-bottom: \textbf{forward substitution}. No searching, no choices -- just a slide.
LU's whole strategy is to trade one hard system for two trivial ones.

\section{Elimination, watched closely}
\label{ch09:elim-watched}

Let us do a single careful elimination and \emph{watch what we throw away}, because the
thing we usually throw away is exactly $\mat{L}$.

Take
\[
  \mat{A} = \begin{bmatrix} 2 & 1 & 1 \\ 4 & 3 & 3 \\ 8 & 7 & 9 \end{bmatrix}.
\]
The first pivot is the $2$ in the top-left corner. To clear the column below it we ask:
what multiple of row $1$ must we subtract from row $2$ to kill the $4$? The answer is
$4/2 = 2$. We call this number the \textbf{multiplier} and name it $\ell_{21}$ (the
multiplier that acts on row $2$ using row $1$):
\[
  \ell_{21} = \frac{4}{2} = 2, \qquad
  R_2 \leftarrow R_2 - 2\,R_1.
\]
Likewise to kill the $8$ in row $3$ we need $\ell_{31} = 8/2 = 4$, and
$R_3 \leftarrow R_3 - 4\,R_1$. After these two operations:
\[
  \begin{bmatrix} 2 & 1 & 1 \\ 4 & 3 & 3 \\ 8 & 7 & 9 \end{bmatrix}
  \;\longrightarrow\;
  \begin{bmatrix} 2 & 1 & 1 \\ 0 & 1 & 1 \\ 0 & 3 & 5 \end{bmatrix}.
\]
The new second pivot is the $1$ now sitting in position $(2,2)$. To clear the $3$ below it,
\[
  \ell_{32} = \frac{3}{1} = 3, \qquad R_3 \leftarrow R_3 - 3\,R_2,
\]
giving
\[
  \begin{bmatrix} 2 & 1 & 1 \\ 0 & 1 & 1 \\ 0 & 3 & 5 \end{bmatrix}
  \;\longrightarrow\;
  \begin{bmatrix} 2 & 1 & 1 \\ 0 & 1 & 1 \\ 0 & 0 & 2 \end{bmatrix} = \mat{U}.
\]
That triangular matrix on the right is $\mat{U}$. Three multipliers did all the work:
\[
  \ell_{21} = 2, \quad \ell_{31} = 4, \quad \ell_{32} = 3.
\]
Now assemble those multipliers into a lower triangular matrix, with
$1$'s on the diagonal and each $\ell_{ij}$ sitting in the exact position $(i,j)$ where it
was used:
\[
  \mat{L} = \begin{bmatrix} 1 & 0 & 0 \\ \ell_{21} & 1 & 0 \\ \ell_{31} & \ell_{32} & 1 \end{bmatrix}
          = \begin{bmatrix} 1 & 0 & 0 \\ 2 & 1 & 0 \\ 4 & 3 & 1 \end{bmatrix}.
\]
Claim: $\mat{A} = \mat{L}\mat{U}$. Let us check by direct multiplication.
\[
  \mat{L}\mat{U}
  = \begin{bmatrix} 1 & 0 & 0 \\ 2 & 1 & 0 \\ 4 & 3 & 1 \end{bmatrix}
    \begin{bmatrix} 2 & 1 & 1 \\ 0 & 1 & 1 \\ 0 & 0 & 2 \end{bmatrix}.
\]
Row $1$ of $\mat{L}$ is $(1,0,0)$, so row $1$ of the product is just row $1$ of $\mat{U}$:
$(2,1,1)$. Row $2$ of $\mat{L}$ is $(2,1,0)$, so row $2$ of the product is
$2\cdot(2,1,1) + 1\cdot(0,1,1) = (4,3,3)$. Row $3$ of $\mat{L}$ is $(4,3,1)$, so row $3$ is
$4\cdot(2,1,1) + 3\cdot(0,1,1) + 1\cdot(0,0,2) = (8,7,9)$. Stacking,
\[
  \mat{L}\mat{U} = \begin{bmatrix} 2 & 1 & 1 \\ 4 & 3 & 3 \\ 8 & 7 & 9 \end{bmatrix} = \mat{A}. \checkmark
\]
It worked for a clean structural reason. The multipliers you used to \emph{destroy} the
lower part of $\mat{A}$ are exactly the numbers that, placed in $\mat{L}$,
\emph{rebuild} it.

\begin{focusbox}
Hold onto three sentences.
\begin{itemize}
  \item $\mat{U}$ is what elimination leaves behind: the upper triangular matrix of pivots.
  \item $\mat{L}$ is the bookkeeping: $1$'s on the diagonal, and the multiplier
        $\ell_{ij}=(\text{entry being cleared})/(\text{pivot})$ in position $(i,j)$.
  \item No minus signs in $\mat{L}$. You \emph{subtracted} $\ell_{ij}$ times a row during
        elimination, but $\mat{L}$ stores the plain multiplier with a $+$ sign.
\end{itemize}
\end{focusbox}

\begin{trapbox}
The single most common LU mistake is putting minus signs into $\mat{L}$. During elimination
you carried out $R_3 \leftarrow R_3 - 3R_2$, so it feels natural to write $-3$ in $\mat{L}$.
\textbf{Do not.} The matrix that \emph{performs} the subtraction is $\mat{E}$ with a $-3$;
but $\mat{L} = \mat{E}\inv$ \emph{undoes} the subtraction, and undoing a subtraction is an
addition. So $\mat{L}$ holds $+3$. We will see exactly why in the next section. Quick sanity
check: if every multiplier is positive (you only ever subtracted positive multiples), then
$\mat{L}$ should be all positive below the diagonal.
\end{trapbox}

\section{Why it works: elimination is left-multiplication by shears}
\label{ch09:why-it-works}

The claim ``$\mat{L}$ stores the multipliers'' deserves a real derivation. The idea behind
it is clean: each elementary row operation is itself a matrix, and undoing it simply
requires reversing the sign of the multiplier.

Recall that performing a row operation on $\mat{A}$ is the same as \emph{left}-multiplying
$\mat{A}$ by an elementary matrix. The operation $R_2 \leftarrow R_2 - 2R_1$ on a
$3\times 3$ matrix is multiplication by
\[
  \mat{E}_{21} = \begin{bmatrix} 1 & 0 & 0 \\ -2 & 1 & 0 \\ 0 & 0 & 1 \end{bmatrix}.
\]
(Read it row by row: new row $2$ $=$ $-2\cdot$ old row $1$ $+1\cdot$ old row $2$. Exactly
the operation.) Our full elimination was three such steps, applied left to right:
\[
  \mat{E}_{32}\,\mat{E}_{31}\,\mat{E}_{21}\,\mat{A} = \mat{U},
\]
with
\[
  \mat{E}_{21} = \begin{bmatrix} 1&0&0\\-2&1&0\\0&0&1 \end{bmatrix},\;
  \mat{E}_{31} = \begin{bmatrix} 1&0&0\\0&1&0\\-4&0&1 \end{bmatrix},\;
  \mat{E}_{32} = \begin{bmatrix} 1&0&0\\0&1&0\\0&-3&1 \end{bmatrix}.
\]

\begin{figure}[H]
\centering
\begin{tikzpicture}[scale=0.88,>=Stealth,every node/.style={font=\small}]
  % State 0: Original matrix A
  \node[draw=sourceblue,thick,rounded corners=2pt,fill=sourceblue!8,inner sep=5pt,align=center] (A0) at (0,0) {
    $\begin{bmatrix}
      \mathbf{2} & 1 & 1 \\
      \color{trapred}4 & 3 & 3 \\
      \color{trapred}8 & 7 & 9
    \end{bmatrix}$\\[4pt]
    \footnotesize $\mat{A}$
  };

  % Transition 1
  \draw[->,thick,sourceblue] (1.3,0) -- (2.7,0)
    node[midway,above,font=\scriptsize]{\color{sourceblue} $\mat{E}_{21}$}
    node[midway,below,font=\tiny]{$R_2 - 2R_1$};

  % State 1: After E_21
  \node[draw=sourceblue,thick,rounded corners=2pt,fill=sourceblue!8,inner sep=5pt,align=center] (A1) at (4.0,0) {
    $\begin{bmatrix}
      \mathbf{2} & 1 & 1 \\
      \color{targetgreen!80!black}\mathbf{0} & \mathbf{1} & 1 \\
      \color{trapred}8 & 7 & 9
    \end{bmatrix}$\\[4pt]
    \footnotesize $\mat{E}_{21}\mat{A}$
  };

  % Transition 2
  \draw[->,thick,sourceblue] (5.3,0) -- (6.7,0)
    node[midway,above,font=\scriptsize]{\color{sourceblue} $\mat{E}_{31}$}
    node[midway,below,font=\tiny]{$R_3 - 4R_1$};

  % State 2: After E_31
  \node[draw=sourceblue,thick,rounded corners=2pt,fill=sourceblue!8,inner sep=5pt,align=center] (A2) at (8.0,0) {
    $\begin{bmatrix}
      \mathbf{2} & 1 & 1 \\
      \color{targetgreen!80!black}\mathbf{0} & \mathbf{1} & 1 \\
      \color{targetgreen!80!black}\mathbf{0} & \color{trapred}3 & 5
    \end{bmatrix}$\\[4pt]
    \footnotesize $\mat{E}_{31}\mat{E}_{21}\mat{A}$
  };

  % Transition 3
  \draw[->,thick,sourceblue] (9.3,0) -- (10.7,0)
    node[midway,above,font=\scriptsize]{\color{sourceblue} $\mat{E}_{32}$}
    node[midway,below,font=\tiny]{$R_3 - 3R_2$};

  % State 3: Final U
  \node[draw=targetgreen!80!black,thick,rounded corners=2pt,fill=targetgreen!10,inner sep=5pt,align=center] (A3) at (12.0,0) {
    $\begin{bmatrix}
      \mathbf{2} & 1 & 1 \\
      \color{targetgreen!80!black}\mathbf{0} & \mathbf{1} & 1 \\
      \color{targetgreen!80!black}\mathbf{0} & \color{targetgreen!80!black}\mathbf{0} & \mathbf{2}
    \end{bmatrix}$\\[4pt]
    \footnotesize $\mat{U}$ (upper triangular)
  };

  % Bottom corridor: Inverting and bundling into L
  \draw[<-,ultra thick,quotientorange] (0.0,-1.7) -- (12.0,-1.7)
    node[midway,above,font=\small]{\bfseries Inverting the shears: $\mat{A} = \mat{L}\mat{U}$ with $\mat{L} = \mat{E}_{21}^{-1}\mat{E}_{31}^{-1}\mat{E}_{32}^{-1}$}
    node[midway,below,font=\scriptsize]{\color{quotientorange!90!black} Multipliers $\ell_{21}=2$, $\ell_{31}=4$, $\ell_{32}=3$ drop cleanly into the lower triangle};
\end{tikzpicture}
\caption{Gaussian elimination as a sequence of elementary shears. Each elementary matrix $\mat{E}_{ij}$ applies a shear that zeroes out the entry below the diagonal in position $(i,j)$ (marked in green). Inverting the sequence in reverse order reconstitutes $\mat{A} = \mat{L}\mat{U}$, where the lower-triangular matrix $\mat{L}$ bundles all the positive multipliers with no mutual interference.}
\label{ch09:fig-shear-sequence}
\end{figure}
Solve for $\mat{A}$ by moving the $\mat{E}$'s to the other side. Each one is invertible
(elimination is reversible -- you can always add back what you subtracted), so
\[
  \mat{A} = \mat{E}_{21}\inv\,\mat{E}_{31}\inv\,\mat{E}_{32}\inv\,\mat{U}.
\]
Now the two gifts. \textbf{Gift one:} inverting a single shear just flips the sign of its
off-diagonal entry, because to undo ``subtract $2$ times row $1$'' you ``add $2$ times row
$1$'':
\[
  \mat{E}_{21}\inv = \begin{bmatrix} 1&0&0\\ +2&1&0\\0&0&1 \end{bmatrix}.
\]
There is the plus sign the trap warned about. \textbf{Gift two:} multiplying these
particular inverse shears together does not mix anything up -- the multipliers simply drop
into their slots with no interaction. Watch:
\[
  \mat{E}_{21}\inv\,\mat{E}_{31}\inv\,\mat{E}_{32}\inv
  = \begin{bmatrix} 1&0&0\\2&1&0\\0&0&1 \end{bmatrix}
    \begin{bmatrix} 1&0&0\\0&1&0\\4&0&1 \end{bmatrix}
    \begin{bmatrix} 1&0&0\\0&1&0\\0&3&1 \end{bmatrix}
  = \begin{bmatrix} 1&0&0\\2&1&0\\4&3&1 \end{bmatrix} = \mat{L}.
\]
So $\mat{A} = \mat{L}\mat{U}$, with $\mat{L}$ exactly the multiplier table. The reason no
cross-terms appear is worth stating once, because it is the whole secret.

\begin{formalbox}
\textbf{Why the multipliers do not interfere.} When you eliminate in the natural order --
column $1$ first (top to bottom), then column $2$, and so on -- every inverse elementary
matrix $\mat{E}_{ij}\inv$ adds a multiple of an \emph{earlier} row $j$ to a \emph{later} row
$i$ (so $i>j$). Multiplying two such matrices, the entry $\ell_{ij}$ of one never gets
touched by the other, because the second factor only modifies rows using columns that the
first factor has already finished with. Concretely, for unit lower triangular elementary
matrices that each have a single off-diagonal entry below the diagonal, the product is just
the matrix containing \emph{all} those entries at once:
\[
  \prod (\mat{E}_{ij}\inv) = \mat{I} + \sum_{i>j} \ell_{ij}\,\vec{e}_i \vec{e}_j\T = \mat{L}.
\]
This is special to the with-the-grain ordering. Eliminate out of order and the clean
``read off the multipliers'' rule breaks.
\end{formalbox}

\begin{walkbox}
Take a walk and chew on this. You did three subtractions to get to $\mat{U}$. To get back
to $\mat{A}$ you must undo them \emph{in reverse order} -- last subtraction undone first.
Yet $\mat{L}$ lists the multipliers in \emph{forward} order with no scrambling. How can the
order not matter for assembling $\mat{L}$ even though it absolutely matters for the row
operations themselves? (Hint: think about which rows each step reads from versus writes to.
A step writes to a row no later step ever reads.)
\end{walkbox}

\begin{strangbox}
\emph{Before Strang:} if you have not yet taken his course, just absorb the slogan
``$\mat{A}=\mat{L}\mat{U}$, where $\mat{U}$ is elimination's output and $\mat{L}$ is the
multipliers.'' That single sentence will make his Lecture on elimination feel like a story
whose ending you already know.

\emph{After Strang:} Strang loves to write elimination as
$\mat{E}\mat{A}=\mat{U}$ and then flip it to $\mat{A}=\mat{E}\inv\mat{U}=\mat{L}\mat{U}$,
emphasizing that $\mat{L}$ is the \emph{nicer} matrix because no minus signs survive. He
also stresses $\mat{LDU}$, which we reach in
Section~\ref{ch09:ldu}: pull the pivots out of $\mat{U}$ into a diagonal $\mat{D}$ so that
both triangular factors have unit diagonals and the asymmetry between them disappears.
\end{strangbox}

\section{Solving \texorpdfstring{$\mat{A}\vec{x}=\vec{b}$}{Ax=b}: two slides instead of one sweep}
\label{ch09:two-slides}

Now we collect the payoff. We want $\mat{A}\vec{x}=\vec{b}$, and we know $\mat{A}=\mat{L}\mat{U}$,
so
\[
  \mat{L}\mat{U}\vec{x} = \vec{b}.
\]
Introduce a helper vector $\vec{c} = \mat{U}\vec{x}$. Then the single problem splits into a
chain of two triangular problems:
\[
  \mat{L}\vec{c} = \vec{b} \quad(\text{forward, solve for }\vec{c}),
  \qquad
  \mat{U}\vec{x} = \vec{c} \quad(\text{backward, solve for }\vec{x}).
\]
First slide down through $\mat{L}$ to find $\vec{c}$, then slide up through $\mat{U}$ to find
$\vec{x}$. Let us do it with our matrix and a concrete right-hand side
$\vec{b} = \colvec{5\\11\\27}$.

\textbf{Forward substitution, $\mat{L}\vec{c}=\vec{b}$.}
\[
  \begin{bmatrix} 1 & 0 & 0 \\ 2 & 1 & 0 \\ 4 & 3 & 1 \end{bmatrix}
  \begin{bmatrix} c_1 \\ c_2 \\ c_3 \end{bmatrix}
  = \begin{bmatrix} 5 \\ 11 \\ 27 \end{bmatrix}.
\]
Row $1$: $c_1 = 5$. Row $2$: $2c_1 + c_2 = 11 \Rightarrow c_2 = 11 - 2(5) = 1$. Row $3$:
$4c_1 + 3c_2 + c_3 = 27 \Rightarrow c_3 = 27 - 20 - 3 = 4$. So
$\vec{c} = \colvec{5\\1\\4}$.

\textbf{Back substitution, $\mat{U}\vec{x}=\vec{c}$.}
\[
  \begin{bmatrix} 2 & 1 & 1 \\ 0 & 1 & 1 \\ 0 & 0 & 2 \end{bmatrix}
  \begin{bmatrix} x_1 \\ x_2 \\ x_3 \end{bmatrix}
  = \begin{bmatrix} 5 \\ 1 \\ 4 \end{bmatrix}.
\]
Row $3$: $2x_3 = 4 \Rightarrow x_3 = 2$. Row $2$: $x_2 + x_3 = 1 \Rightarrow x_2 = 1 - 2 = -1$.
Row $1$: $2x_1 + x_2 + x_3 = 5 \Rightarrow 2x_1 = 5 - (-1) - 2 = 4 \Rightarrow x_1 = 2$. So
\[
  \vec{x} = \begin{bmatrix} 2 \\ -1 \\ 2 \end{bmatrix}.
\]
Check against the original: $\mat{A}\vec{x} = \colvec{2\\4\\8}\!(2) + \colvec{1\\3\\7}\!(-1)
+ \colvec{1\\3\\9}\!(2) = \colvec{4-1+2 \\ 8-3+6 \\ 16-7+18} = \colvec{5\\11\\27}
= \vec{b}$. \checkmark

\begin{figure}[H]
\centering
\begin{tikzpicture}[>=Stealth,node distance=8mm]
  \tikzstyle{box}=[draw,rounded corners,minimum width=20mm,minimum height=9mm,align=center]
  \node[box,fill=blue!5] (b) {$\vec{b}$};
  \node[box,fill=green!6,right=22mm of b] (c) {$\vec{c}$};
  \node[box,fill=orange!6,right=22mm of c] (x) {$\vec{x}$};
  \draw[->,thick] (b) -- node[above,font=\small]{$\mat{L}\vec{c}=\vec{b}$} node[below,font=\small\itshape]{forward} (c);
  \draw[->,thick] (c) -- node[above,font=\small]{$\mat{U}\vec{x}=\vec{c}$} node[below,font=\small\itshape]{backward} (x);
\end{tikzpicture}
\caption{Solving with a factorization is a two-stop journey. Slide forward through the lower
triangular $\mat{L}$ to convert $\vec{b}$ into the intermediate $\vec{c}$, then slide
backward through the upper triangular $\mat{U}$ to land on the answer $\vec{x}$. Both slides
are mechanical; neither involves any elimination.}
\label{ch09:fig-twoslides}
\end{figure}

The crucial point is what happened to the cost. Building $\mat{L}$ and $\mat{U}$ took the
full elimination effort -- once. But the two substitutions are cheap, and they are
\emph{all you need} for a brand-new $\vec{b}$. We will quantify ``cheap'' soon, but the
shape of the win is already visible: \textbf{factor once, substitute many times}.

\begin{labbox}
By hand, factor
\[
  \mat{A} = \begin{bmatrix} 3 & 1 \\ 6 & 5 \end{bmatrix}
\]
into $\mat{L}\mat{U}$. (Multiplier $\ell_{21}=6/3=2$; subtract $2$ times row $1$ from row
$2$.) You should get $\mat{L}=\begin{bmatrix}1&0\\2&1\end{bmatrix}$,
$\mat{U}=\begin{bmatrix}3&1\\0&3\end{bmatrix}$. Then solve $\mat{A}\vec{x}=\colvec{4\\11}$ by
forward-then-back substitution and confirm $\vec{x}=\colvec{1\\1}$ (you should get the
intermediate $\vec{c}=\colvec{4\\3}$). Finally multiply
$\mat{L}\mat{U}$ back out and watch $\mat{A}$ reappear.
\end{labbox}

\section{Many right-hand sides: the real reason to factor}
\label{ch09:many-rhs}

Here is the situation that makes LU indispensable. Suppose you must solve
\[
  \mat{A}\vec{x}_1 = \vec{b}_1,\quad \mat{A}\vec{x}_2 = \vec{b}_2,\quad \dots,\quad
  \mat{A}\vec{x}_k = \vec{b}_k,
\]
with the \emph{same} matrix $\mat{A}$ but many different right-hand sides. This is the
common case, not the exception. In circuit simulation $\mat{A}$ is the network and each
$\vec{b}$ is a different set of driving voltages. In a Newton solver for a nonlinear
problem, $\mat{A}$ is a Jacobian reused across iterations. In finite-element structural
analysis $\mat{A}$ is the stiffness matrix and each $\vec{b}$ is a different load case.

The naive approach re-runs elimination $k$ times: a recompute. The LU approach factors
$\mat{A}=\mat{L}\mat{U}$ \emph{once} and then does $k$ cheap forward/back pairs. Since
elimination is far more expensive than substitution, the savings grow without bound as $k$
grows.

There is a clean way to package ``all the right-hand sides at once.'' Put them as columns of
a matrix $\mat{B}=[\vec{b}_1\ \vec{b}_2\ \cdots\ \vec{b}_k]$ and solve $\mat{A}\mat{X}=\mat{B}$
column by column, reusing the one factorization. A particularly important special case is
$\mat{B}=\mat{I}$: solving $\mat{A}\mat{X}=\mat{I}$ gives $\mat{X}=\mat{A}\inv$. So the
matrix inverse is just LU applied to the $n$ right-hand sides $\vec{e}_1,\dots,\vec{e}_n$.

\begin{trapbox}
This last fact is also a warning. To \emph{solve} $\mat{A}\vec{x}=\vec{b}$ you almost never
want to form $\mat{A}\inv$ first and multiply. Computing $\mat{A}\inv$ is like doing $n$ full
solves and then throwing in an extra matrix-vector multiply at the end -- more work, and
numerically worse. The professional habit is: factor once, then forward/back-substitute. If
someone's code contains \texttt{inv(A)*b}, they are doing it the slow, fragile way; the
right call is a triangular solve. ``Never invert; always factor.''
\end{trapbox}

\begin{applicationbox}
\textbf{Real-time engineering and the power grid.} An electrical power grid is modeled by a
large sparse system $\mat{A}\vec{x}=\vec{b}$, where $\mat{A}$ encodes the network admittances
(which buses connect to which, and how strongly) and $\vec{b}$ encodes the loads and
generation at this instant. The grid's wiring -- the matrix $\mat{A}$ -- changes only when a
line trips or a switch flips, which is rare. But the loads $\vec{b}$ change every few
seconds as demand fluctuates. So operators factor $\mat{A}=\mat{L}\mat{U}$ once (for sparse
matrices a specialized sparse LU keeps $\mat{L}$ and $\mat{U}$ almost as sparse as $\mat{A}$)
and then re-solve thousands of times a day with fresh $\vec{b}$'s by cheap substitution.
This same ``factor once, substitute often'' pattern drives implicit time-stepping in fluid
and climate simulation, interior-point optimization, real-time robotics control, and the
linear solves buried inside training large statistical models. LU is one of the most-executed
numerical algorithms in human history precisely because of this reuse.
\end{applicationbox}

\section{Row exchanges: when a pivot is missing}
\label{ch09:pivoting}

Everything above assumed elimination glided through without ever hitting a zero where it
needed a pivot. Real matrices are not so polite. Consider
\[
  \mat{A} = \begin{bmatrix} 0 & 1 \\ 2 & 3 \end{bmatrix}.
\]
The would-be first pivot is the top-left $0$. You cannot divide by zero to form a
multiplier, so elimination stalls -- not because the matrix is bad (it is perfectly
invertible) but because the rows are simply in an inconvenient order. The cure is obvious:
swap the two rows.
\[
  \begin{bmatrix} 0 & 1 \\ 2 & 3 \end{bmatrix}
  \xrightarrow{\ R_1 \leftrightarrow R_2\ }
  \begin{bmatrix} 2 & 3 \\ 0 & 1 \end{bmatrix}.
\]
Now there is a healthy pivot $2$, and the matrix is already upper triangular
(Figure~\ref{ch09:fig-swap}). A row swap is
itself a matrix -- a \textbf{permutation matrix} $\mat{P}$, which is the identity with its
rows reordered:
\[
  \mat{P} = \begin{bmatrix} 0 & 1 \\ 1 & 0 \end{bmatrix},
  \qquad
  \mat{P}\mat{A} = \begin{bmatrix} 0&1\\1&0 \end{bmatrix}\begin{bmatrix} 0&1\\2&3 \end{bmatrix}
  = \begin{bmatrix} 2 & 3 \\ 0 & 1 \end{bmatrix}.
\]
Left-multiplying by $\mat{P}$ reorders the rows of whatever it hits. Permutation matrices
are wonderfully simple: each row and each column has exactly one $1$ and the rest $0$, they
just shuffle, and their inverse is their transpose ($\mat{P}\inv = \mat{P}\T$, since undoing
a shuffle is a shuffle).

\begin{figure}[H]
\centering
\begin{tikzpicture}[>=Stealth,every node/.style={font=\small}]
  \matrix (m) [matrix of math nodes, left delimiter={[}, right delimiter={]},
               row sep=2pt, column sep=10pt,
               nodes={minimum width=6mm}] {
    0 & 1 \\
    2 & 3 \\
  };
  \node at (3.0,0) {$\xrightarrow{\;R_1\leftrightarrow R_2\;}$};
  \matrix (n) [matrix of math nodes, left delimiter={[}, right delimiter={]},
               row sep=2pt, column sep=10pt, xshift=5.6cm,
               nodes={minimum width=6mm}] {
    2 & 3 \\
    0 & 1 \\
  };
  % arrows showing the swap of rows
  \draw[->,vecred,thick] (m-1-2.east) to[bend left=30] node[above,font=\scriptsize,vecred] {top row $\downarrow$} ($(n-2-1.west)+(-0.15,0)$);
  \draw[->,vecblue,thick] (m-2-2.east) to[bend right=30] node[below,font=\scriptsize,vecblue] {bottom row $\uparrow$} ($(n-1-1.west)+(-0.15,0)$);
  \node[font=\scriptsize,gray] at (n.south)[below=4mm] {pivot $2\neq0$ now in place};
\end{tikzpicture}
\caption{A row exchange is the whole content of a permutation matrix $\mat{P}$. The dead
zero in the pivot slot (top-left) is swapped out for a live entry by interchanging the two
rows; left-multiplying by $\mat{P}=\protect\begin{bmatrix}0&1\\1&0\protect\end{bmatrix}$ does
exactly this. After the swap the matrix is already upper triangular, so elimination is done.}
\label{ch09:fig-swap}
\end{figure}

The fully general statement of the factorization, valid for \emph{every} invertible matrix,
builds the permutation in:

\begin{formalbox}
\textbf{LU with partial pivoting.} For any invertible square matrix $\mat{A}$ there is a
permutation matrix $\mat{P}$, a unit lower triangular $\mat{L}$ (ones on the diagonal), and
an upper triangular $\mat{U}$ such that
\[
  \mat{P}\mat{A} = \mat{L}\mat{U}.
\]
Read it as: \emph{if you reorder the rows of $\mat{A}$ in the right way up front, the
reordered matrix factors cleanly into $\mat{L}\mat{U}$.} The permutation $\mat{P}$ records
all the swaps elimination chose to make.
\end{formalbox}

To solve $\mat{A}\vec{x}=\vec{b}$ from $\mat{P}\mat{A}=\mat{L}\mat{U}$, multiply the system
by $\mat{P}$ first: $\mat{P}\mat{A}\vec{x}=\mat{P}\vec{b}$, i.e.
$\mat{L}\mat{U}\vec{x}=\mat{P}\vec{b}$. So you permute the right-hand side the same way, then
do the usual forward/back substitution. Permuting $\vec{b}$ is free -- it is just reading its
entries in a different order.

\subsection{Pivoting is not only for zeros -- it is for stability}
\label{ch09:pivot-stability}

Here is the deeper reason real software \emph{always} pivots, even when no pivot is exactly
zero. A \emph{tiny} pivot is nearly as bad as a zero pivot. Suppose
\[
  \mat{A} = \begin{bmatrix} 0.0001 & 1 \\ 1 & 1 \end{bmatrix}.
\]
Use the tiny $0.0001$ as a pivot: the multiplier is $\ell_{21}=1/0.0001 = 10000$, and
clearing row $2$ gives $u_{22} = 1 - 10000\cdot 1 = -9999$. That huge multiplier amplifies
any rounding error in the data by a factor of $10^4$. The fix is \textbf{partial pivoting}:
before using a column, swap the row with the \emph{largest} available entry into the pivot
position, so every multiplier satisfies $|\ell_{ij}|\le 1$. Swapping rows here puts the $1$
on top, the multiplier becomes a tame $0.0001$, and the computation stays accurate. So
$\mat{P}$ is doing double duty: it rescues you from exact zeros, and it keeps the arithmetic
numerically stable. We return to this in the numerical chapter; for now, just know that
$\mat{P}\mat{A}=\mat{L}\mat{U}$ is what your computer actually does.

\begin{example}[A full $\mat{P}\mat{A}=\mat{L}\mat{U}$]
\label{ch09:ex-palu}
Factor $\mat{A}=\begin{bmatrix} 1 & 2 \\ 3 & 4 \end{bmatrix}$ with partial pivoting. Column
$1$ has entries $1$ and $3$; the larger magnitude is $3$, in row $2$, so swap first:
\[
  \mat{P} = \begin{bmatrix} 0&1\\1&0 \end{bmatrix},\quad
  \mat{P}\mat{A} = \begin{bmatrix} 3 & 4 \\ 1 & 2 \end{bmatrix}.
\]
Now eliminate: multiplier $\ell_{21}=1/3$, and
$R_2 \leftarrow R_2 - \tfrac13 R_1$ gives $u_{22} = 2 - \tfrac13\cdot 4 = \tfrac23$. Hence
\[
  \mat{L} = \begin{bmatrix} 1 & 0 \\ \tfrac13 & 1 \end{bmatrix}, \qquad
  \mat{U} = \begin{bmatrix} 3 & 4 \\ 0 & \tfrac23 \end{bmatrix}.
\]
Verify: $\mat{L}\mat{U} = \begin{bmatrix} 3 & 4 \\ 1 & \tfrac43 + \tfrac23 \end{bmatrix}
= \begin{bmatrix} 3 & 4 \\ 1 & 2 \end{bmatrix} = \mat{P}\mat{A}$. \checkmark
\end{example}

\begin{connectionbox}
\textbf{The sign of a permutation, and the determinant.} A permutation matrix is the matrix
face of a permutation of $\{1,\dots,n\}$, the objects group theory calls $S_n$. Every
permutation is a product of swaps, and although the number of swaps is not unique, its
\emph{parity} (even or odd) is -- a fixed fact about the permutation. That parity is the
\emph{sign}, $\det\mat{P}=\pm 1$. This links straight to determinants:
$\mat{P}\mat{A}=\mat{L}\mat{U}$ gives
\[
  \det\mat{A} = (\det\mat{P})\,(\det\mat{L})\,(\det\mat{U})
              = (\pm 1)\cdot 1 \cdot (u_{11}u_{22}\cdots u_{nn}),
\]
because $\det$ of a triangular matrix is the product of its diagonal, and $\mat{L}$ has unit
diagonal. So the determinant is just $\pm$ the product of the pivots -- which is exactly how
numerical libraries compute determinants: factor, multiply the pivots, fix the sign from the
swaps. Elimination, permutations, and determinants are three views of one object.
\end{connectionbox}

\section{Pulling out the pivots: \texorpdfstring{$\mat{A}=\mat{L}\mat{D}\mat{U}$}{A=LDU}}
\label{ch09:ldu}

There is a mild asymmetry in $\mat{A}=\mat{L}\mat{U}$ that bothers people, and fixing it
reveals structure. In our example
\[
  \mat{L} = \begin{bmatrix} 1&0&0\\2&1&0\\4&3&1 \end{bmatrix}, \qquad
  \mat{U} = \begin{bmatrix} 2&1&1\\0&1&1\\0&0&2 \end{bmatrix},
\]
the $\mat{L}$ has $1$'s on its diagonal but $\mat{U}$ carries the pivots $2,1,2$ on its
diagonal. Factor those pivots out of $\mat{U}$ into a diagonal matrix $\mat{D}$, leaving a
$\mat{U}$ with unit diagonal too:
\[
  \mat{U} = \underbrace{\begin{bmatrix} 2&0&0\\0&1&0\\0&0&2 \end{bmatrix}}_{\mat{D}}
            \underbrace{\begin{bmatrix} 1 & \tfrac12 & \tfrac12 \\ 0 & 1 & 1 \\ 0 & 0 & 1 \end{bmatrix}}_{\mat{U}'},
\]
where each row of $\mat{U}$ was divided by its pivot. So
\[
  \mat{A} = \mat{L}\,\mat{D}\,\mat{U}', \qquad
  \mat{D}=\diag(2,1,2),
\]
with $\mat{L}$ and $\mat{U}'$ both unit-triangular and all the scaling living in $\mat{D}$.
This $\mat{L}\mat{D}\mat{U}$ form is the symmetric-looking sibling of $\mat{L}\mat{U}$, and it
pays off for symmetric matrices: if $\mat{A}=\mat{A}\T$ then $\mat{U}'=\mat{L}\T$ and the
factorization collapses to $\mat{A}=\mat{L}\mat{D}\mat{L}\T$. When additionally all pivots in
$\mat{D}$ are positive you can split $\mat{D}=\sqrt{\mat{D}}\sqrt{\mat{D}}$ and absorb the
roots, landing on $\mat{A}=\mat{R}\T\mat{R}$ -- the celebrated \textbf{Cholesky
factorization}, the symmetric, half-the-work specialization of LU that powers least squares
and Gaussian statistics. We will meet it again with positive definite matrices.

\begin{aestheticbox}
What makes LU lovely is that it turns a \emph{process} into an \emph{object}. Elimination is
something you \emph{do} -- a sequence of moves unfolding in time. The factorization
$\mat{A}=\mat{L}\mat{U}$ freezes that entire process into two static matrices you can store,
inspect, multiply, and reuse. The fleeting becomes permanent. And the same matrix sprouts a
family: drop minus signs and you get $\mat{L}$ from elimination; pull out pivots and you get
$\mat{L}\mat{D}\mat{U}$; impose symmetry and it folds to $\mat{L}\mat{D}\mat{L}\T$; demand
positivity and it crystallizes into Cholesky. One idea, refracted through different
structures, lights up the whole subject.
\end{aestheticbox}

\section{The cost of elimination, counted exactly}
\label{ch09:cost}

We have been saying ``elimination is expensive, substitution is cheap.'' Let us make that
quantitative, because the numbers are the entire justification for the LU strategy.

Count multiply-and-subtract operations (``flops''). To clear column $1$ of an $n\times n$
matrix you process rows $2$ through $n$ -- that is $n-1$ rows -- and for each you compute one
multiplier and update about $n$ entries, so roughly $(n-1)\cdot n$ operations. Clearing
column $2$ works on the remaining $(n-1)\times(n-1)$ block, costing about $(n-2)(n-1)$. In
general, eliminating column $k$ costs about $(n-k)(n-k+1)$. Summing,
\[
  \text{elimination cost} \;\approx\; \sum_{k=1}^{n-1} (n-k)^2 \;\approx\; \sum_{j=1}^{n-1} j^2
  \;\approx\; \frac{n^3}{3}.
\]
The leading term is $\tfrac13 n^3$. \textbf{Cubic.} Double the size of the matrix and
elimination costs eight times as much.

Now count a triangular solve. Back substitution on $\mat{U}$: the last unknown is free, the
next costs $1$ multiply, the next $2$, and so on up to about $n$, giving
\[
  \sum_{k=1}^{n} k \approx \frac{n^2}{2}.
\]
Forward substitution on $\mat{L}$ costs the same $\tfrac12 n^2$. So one full solve from a
finished factorization is about $n^2$ operations. \textbf{Quadratic.}

\begin{figure}[H]
\centering
\begin{tikzpicture}
\begin{axis}[
  width=10.5cm, height=6.4cm,
  xlabel={matrix size $n$}, ylabel={operations},
  domain=1:60, samples=80,
  legend pos=north west, legend cell align=left,
  ymin=0, ymax=72000,
  grid=both, grid style={gray!18},
  every axis plot/.append style={very thick},
]
\addplot[vecred]   {x^3/3};   \addlegendentry{$\tfrac13 n^3$ \ (factor once)}
\addplot[vecblue]  {x^2};     \addlegendentry{$n^2$ \ (each extra solve)}
\end{axis}
\end{tikzpicture}
\caption{The two costs that justify LU. Building the factorization is cubic ($\sim\tfrac13
n^3$); each subsequent solve with a new right-hand side is only quadratic ($\sim n^2$). For
large $n$ the gap is enormous, so reusing the factorization across many right-hand sides is
where the savings live.}
\label{ch09:fig-cost}
\end{figure}

The ratio of these tells the whole story. The expensive factorization costs about
$\tfrac13 n^3$; each additional solve costs about $n^2$. Their ratio is
\[
  \frac{\tfrac13 n^3}{n^2} = \frac{n}{3}.
\]
So one factorization ``is worth'' roughly $n/3$ solves. For a modest $n=300$, the factor
step costs as much as a hundred substitution steps. If you have a hundred right-hand sides
and you re-eliminate each time, you do a hundred cubic operations; if you factor once and
substitute, you do one cubic plus a hundred quadratics -- essentially the cost of a
\emph{single} naive solve. That is the difference between a job finishing in a second and
finishing in a minute.

\begin{example}[Counting the win]
\label{ch09:ex-count}
Take $n=1000$ and $k=500$ right-hand sides. Re-eliminating each time costs about
$500 \times \tfrac13(1000)^3 \approx 1.7\times 10^{11}$ operations. Factoring once and
substituting costs about $\tfrac13(1000)^3 + 500\times(1000)^2 \approx 3.3\times 10^8 +
5\times 10^8 \approx 8.3\times 10^8$ operations. The LU strategy is roughly $200$ times
faster here -- and the advantage only widens as $k$ grows, because every extra right-hand
side is nearly free.
\end{example}

\begin{focusbox}
The economics of LU in one breath: \textbf{factoring is cubic, solving is quadratic.} You
pay the cubic price once and then enjoy quadratic solves forever after. Whenever you see the
same matrix with changing right-hand sides -- and you will see it constantly -- factor it.
\end{focusbox}

\section{When LU fails, and when it does not}
\label{ch09:existence}

A short, careful accounting of the boundaries.

\textbf{Pure $\mat{L}\mat{U}$ (no permutation) exists} exactly when elimination never needs a
swap, i.e. when every pivot turns out nonzero without reordering. A clean sufficient
condition: if every \emph{leading principal submatrix} of $\mat{A}$ (the top-left
$1\times 1$, $2\times 2$, \dots, $k\times k$ blocks) is invertible, then $\mat{A}=\mat{L}\mat{U}$
exists and is unique. Uniqueness is easy to believe: if $\mat{L}_1\mat{U}_1=\mat{L}_2\mat{U}_2$
then $\mat{L}_2\inv\mat{L}_1=\mat{U}_2\mat{U}_1\inv$; the left side is unit lower triangular,
the right side is upper triangular, and the only matrix that is both is the identity -- so
$\mat{L}_1=\mat{L}_2$ and $\mat{U}_1=\mat{U}_2$.

\textbf{$\mat{P}\mat{A}=\mat{L}\mat{U}$ always exists} for an invertible $\mat{A}$. There is
always \emph{some} reordering of the rows that makes pivots appear, because invertibility
guarantees a nonzero entry to promote in each column as you go.

\textbf{Singular matrices} can still be reduced to an upper triangular $\mat{U}$, but $\mat{U}$
will have a zero pivot somewhere, signaling that $\mat{A}\vec{x}=\vec{b}$ either has no
solution or infinitely many. That is a story for the chapters on rank and the four
subspaces; LU's clean two-slide solve assumes the pivots are all there.

\begin{remark}
A subtle point worth flagging: the matrix $\begin{bmatrix}0&1\\2&3\end{bmatrix}$ from
Section~\ref{ch09:pivoting} is invertible yet has \emph{no} plain $\mat{L}\mat{U}$
factorization -- its leading $1\times 1$ block is the zero matrix $[0]$, which is not
invertible, so the leading-submatrix test fails at the very first step. This is precisely
the case the permutation rescues. ``No $\mat{L}\mat{U}$'' is a statement about row order, not
about the matrix being defective.
\end{remark}

\begin{strangbox}
\emph{After Strang:} when Strang asks ``which matrices have an $\mat{L}\mat{U}$ without row
exchanges?'', the answer is the leading-submatrix condition above. He frames the general
case as $\mat{P}\mat{A}=\mat{L}\mat{U}$ and emphasizes that this is the version your software
performs. He also connects $\mat{A}=\mat{L}\mat{U}$ to the \emph{four subspaces} later: the
pivots in $\mat{U}$ reveal the rank, and the columns of $\mat{L}$ relate to the column space.
Keep this chapter handy -- it is the computational engine those structural chapters run on.
\end{strangbox}

\section{What to carry forward}
\label{ch09:summary}

\begin{itemize}
  \item \textbf{Elimination is factorization.} Running Gaussian elimination on $\mat{A}$
        secretly produces $\mat{A}=\mat{L}\mat{U}$: $\mat{U}$ is the upper triangular result,
        and $\mat{L}$ collects the multipliers (with a plus sign) below a unit diagonal.
  \item \textbf{Solve in two slides.} $\mat{A}\vec{x}=\vec{b}$ becomes $\mat{L}\vec{c}=\vec{b}$
        (forward) then $\mat{U}\vec{x}=\vec{c}$ (backward). Each slide is mechanical.
  \item \textbf{Factor once, reuse forever.} New right-hand sides need only the two cheap
        substitutions. This is the entire reason LU dominates practical computing.
  \item \textbf{Row exchanges give $\mat{P}\mat{A}=\mat{L}\mat{U}$.} Pivoting handles missing
        pivots and prevents numerical errors from blowing up by keeping multipliers
        $\le 1$ in magnitude.
  \item \textbf{Cost.} Factoring is $\sim\tfrac13 n^3$ (cubic); each solve is $\sim n^2$
        (quadratic). One factorization is worth about $n/3$ solves.
  \item \textbf{Never invert; always factor.} To solve a system, do not compute $\mat{A}\inv$;
        factor and substitute. It is faster and more accurate.
\end{itemize}

You now own the algorithm that quietly sits underneath ``solve this linear system'' in every
serious numerical library. The chapters ahead build the \emph{structure} of linear algebra --
subspaces, rank, orthogonality -- but when those theories need to actually compute, this is
the engine they reach for.

\begin{historybox}
The modern formulation of elimination as matrix factorization ($\mat{A} = \mat{L}\mat{U}$) was established
by Alan Turing in his monumental 1948 paper, \emph{Rounding-Off Errors in Matrix Processes}. Turing, working
on the design of the ACE (Automatic Computing Engine) at the National Physical Laboratory in Britain, wanted
to know whether digital computers could solve systems of one hundred linear equations without roundoff error
destroying the results. In this single paper, Turing introduced the matrix condition number, derived backward
error bounds, and proved that Gaussian elimination is mathematically equivalent to factoring a matrix into
lower and upper triangular factors: $\mat{A} = \mat{L}\mat{U}$.

An important precursor to Turing's work came from military geodesy. In 1910, the French artillery officer and
surveyor André-Louis Cholesky developed an ingenious method for solving the symmetric, positive-definite normal
equations arising in map triangulation. Cholesky was killed in action in August 1918 during the closing months
of World War~I, but his fellow officer Commandant Benoit published Cholesky's manuscript posthumously in 1924.
The Cholesky factorization ($\mat{A} = \mat{L}\mat{L}\T$) remains the gold standard for symmetric
positive-definite systems, requiring only half the arithmetic and half the storage of standard $LU$.

In the early computer era of the 1940s and 1950s, computer memory was measured in mere hundreds of words.
Storing a second $n\times n$ matrix for elimination was out of the question. Turing and subsequent numerical
programmers realized an astonishing fact: the diagonal of $\mat{L}$ is entirely $1$'s, so it does not need to
be stored at all! As elimination wipes out each zero below the pivot in column $j$, the multiplier $\ell_{ij}$
can be written directly into the memory slot where that zero was just created. At the end of elimination,
the original array $\mat{A}$ holds both $\mat{L}$ (below the diagonal) and $\mat{U}$ (on and above the diagonal)
in place with zero additional memory overhead!

The $LU$ factorization transformed Gaussian elimination from a disposable,
destructive scratchpad procedure into a permanent, reusable mathematical asset. By factoring $\mat{A} = \mat{L}\mat{U}$
once at cost $\tfrac{1}{3}n^3$, any subsequent system $\mat{A}\vec{x} = \vec{b}$ with the same matrix can be solved
via forward substitution ($\mat{L}\vec{c}=\vec{b}$) and backward substitution ($\mat{U}\vec{x}=\vec{c}$) in only
$n^2$ operations -- a massive computational saving. Yet standard $LU$ factorization requires square, full-rank
matrices and unpermuted pivots. It cannot solve overdetermined systems where there are more equations than unknowns
($m > n$, as in experimental data fitting), nor can it handle singular or rectangular operators without pivoting.
Moreover, the triangular factors $\mat{L}$ and $\mat{U}$ are not orthogonal; they can distort angles and stretch
lengths. Extending elimination to general matrices requires permutations and transposes: Chapter~\ref{ch10:def-transpose}
introduces permutations and transposes, completing the full decomposition $\mat{P}\mat{A} = \mat{L}\mat{U}$ and
opening the door to symmetric matrices, orthogonality, and the geometry of subspaces in Part~3.
\end{historybox}

\section{Exercises}
\label{ch09:exercises}

\begin{exercise}[Warmup factorization]
\label{ch09:ex-warmup}
Factor $\mat{A}=\begin{bmatrix} 1 & 2 \\ 4 & 9 \end{bmatrix}$ into $\mat{L}\mat{U}$ by hand.
State the multiplier you used, and verify your answer by multiplying $\mat{L}\mat{U}$ back
out.
\end{exercise}

\begin{exercise}[Two slides]
\label{ch09:ex-twoslides}
Using your factorization from the previous exercise, solve $\mat{A}\vec{x}=\colvec{5\\21}$ by
forward substitution ($\mat{L}\vec{c}=\vec{b}$) followed by back substitution
($\mat{U}\vec{x}=\vec{c}$). Show $\vec{c}$ and $\vec{x}$.
\end{exercise}

\begin{exercise}[A $3\times3$ by hand]
\label{ch09:ex-3x3}
Factor
\[
  \mat{A} = \begin{bmatrix} 1 & 1 & 1 \\ 2 & 3 & 5 \\ 4 & 6 & 8 \end{bmatrix}
\]
into $\mat{L}\mat{U}$. Record all three multipliers $\ell_{21},\ell_{31},\ell_{32}$, and
confirm $\mat{L}\mat{U}=\mat{A}$.
\end{exercise}

\begin{exercise}[No minus signs]
\label{ch09:ex-signs}
A student factors $\mat{A}=\begin{bmatrix}2&4\\6&3\end{bmatrix}$ and writes
$\mat{L}=\begin{bmatrix}1&0\\-3&1\end{bmatrix}$. Without recomputing from scratch, explain
how you know the sign is wrong, give the correct $\mat{L}$, and state the general rule for the
sign of the entries of $\mat{L}$.
\end{exercise}

\begin{exercise}[Permutation needed]
\label{ch09:ex-perm}
Show that $\mat{A}=\begin{bmatrix} 0 & 2 \\ 3 & 1 \end{bmatrix}$ has no plain $\mat{L}\mat{U}$
factorization. Then find a permutation $\mat{P}$ so that $\mat{P}\mat{A}=\mat{L}\mat{U}$, and
give $\mat{P}$, $\mat{L}$, and $\mat{U}$ explicitly.
\end{exercise}

\begin{exercise}[Partial pivoting for stability]
\label{ch09:ex-pivot-stability}
Consider $\mat{A}=\begin{bmatrix} 0.001 & 1 \\ 1 & 2 \end{bmatrix}$. (a) Factor it
\emph{without} pivoting and report the multiplier $\ell_{21}$ and the entry $u_{22}$.
(b) Now factor $\mat{P}\mat{A}$ with the rows swapped (partial pivoting) and report the new
multiplier. (c) In one sentence, say why the second factorization is the one a numerical
library prefers.
\end{exercise}

\begin{exercise}[Pull out the pivots]
\label{ch09:ex-ldu}
For $\mat{A}=\begin{bmatrix} 2 & 6 \\ 6 & 20 \end{bmatrix}$, find $\mat{L}\mat{U}$, then pull
the pivots out to write $\mat{A}=\mat{L}\mat{D}\mat{U}'$ with both triangular factors having
unit diagonal. Verify that $\mat{A}$ is symmetric and that, consequently, $\mat{U}'=\mat{L}\T$.
\end{exercise}

\begin{exercise}[Determinant from pivots]
\label{ch09:ex-det}
Using the factorization from Exercise~\ref{ch09:ex-3x3}, compute $\det\mat{A}$ as the product
of the pivots (the diagonal of $\mat{U}$). Then state how the answer would change if one row
exchange had been required.
\end{exercise}

\begin{exercise}[Counting the cost]
\label{ch09:ex-cost}
You must solve $\mat{A}\vec{x}=\vec{b}$ for an $n\times n$ matrix with $n=200$, across
$k=50$ different right-hand sides. (a) Estimate the total operation count if you re-run
elimination ($\sim\tfrac13 n^3$ each) every time. (b) Estimate it if you factor once
($\sim\tfrac13 n^3$) and then do $k$ substitutions ($\sim n^2$ each). (c) What is the
approximate speedup factor?
\end{exercise}

\begin{exercise}[Why not invert? (short proof/argument)]
\label{ch09:ex-noinvert}
Explain, in terms of operation counts, why solving $\mat{A}\vec{x}=\vec{b}$ by first
computing $\mat{A}\inv$ and forming $\mat{A}\inv\vec{b}$ is more expensive than factoring and
substituting. (Hint: how many right-hand sides does computing $\mat{A}\inv$ secretly solve,
and how does that compare to the one right-hand side you actually care about?)
\end{exercise}

\begin{exercise}[Uniqueness, a clean proof]
\label{ch09:ex-uniqueness}
Suppose a matrix has two unit-lower/upper factorizations
$\mat{L}_1\mat{U}_1=\mat{L}_2\mat{U}_2$, where $\mat{L}_1,\mat{L}_2$ are unit lower triangular
and $\mat{U}_1,\mat{U}_2$ are upper triangular with nonzero diagonals. Prove $\mat{L}_1=\mat{L}_2$
and $\mat{U}_1=\mat{U}_2$. (Hint: rearrange to $\mat{L}_2\inv\mat{L}_1=\mat{U}_2\mat{U}_1\inv$
and argue about what kind of matrix each side must be.)
\end{exercise}

\begin{exercise}[Block thinking, a challenge]
\label{ch09:ex-block}
Let $\mat{A}=\begin{bmatrix} \mat{A}_{11} & \mat{A}_{12} \\ \mat{A}_{21} & \mat{A}_{22}
\end{bmatrix}$ in $2\times 2$ block form with $\mat{A}_{11}$ invertible. Show that
\[
  \mat{A} = \begin{bmatrix} \mat{I} & \mat{0} \\ \mat{A}_{21}\mat{A}_{11}\inv & \mat{I} \end{bmatrix}
            \begin{bmatrix} \mat{A}_{11} & \mat{A}_{12} \\ \mat{0} & \mat{S} \end{bmatrix},
  \qquad \mat{S} = \mat{A}_{22} - \mat{A}_{21}\mat{A}_{11}\inv\mat{A}_{12}.
\]
This is one step of block elimination; $\mat{S}$ is the \emph{Schur complement}. Identify
which factor plays the role of $\mat{L}$ and which plays $\mat{U}$, and explain how this
generalizes the scalar multiplier $\ell_{21}=a_{21}/a_{11}$.
\end{exercise}
